\section{Cubic fourfolds of admissible discriminant}
\label{cub:section}

We now combine the two birational calculations.  The first records an
integral transcendental contribution in a smooth factorization; the second
rules out that contribution along a Mori factorization, uniformly as the
degree tends to infinity.
On a cubic $X$, write $H$ for the hyperplane divisor and
$h=c_1(\mathcal O_X(1))$ for its cohomology class.

\begin{proposition}[An irrationality criterion]
\label{cub:criterion}
There is an integer $d_0$ with the following property.  Let $X$ be a smooth
complex cubic fourfold, and let
\[
 A_X=H^4(X,\Z)\cap H^{2,2}(X),
 \qquad T_X=A_X^\perp\subset H^4(X,\Z).
\]
Suppose that
\[
 \operatorname{rank}T_X=21,\qquad
 d=|\det T_X|>d_0,\qquad
 \operatorname{End}_{\mathrm{Hdg}}(T_X\otimes\Q)=\Q.
\]
Then $X$ is not rational.
\end{proposition}

\begin{proof}
Choose $d_0>2$ at least as large as the uniform bound in
Theorem~\ref{sf:vanishing}.  Its uniformity allows the collection of K3
surfaces to be restricted to any specified integral transcendental
Hodge-isometry class.

Suppose that $f\colon\mathbb P^4\dashrightarrow X$ is birational.  Take as
test collection the Picard-number-one K3 surfaces of degree $d$ whose
transcendental Hodge lattice, after the shift of Hodge types and reversal
of the cup pairing, is isometric to $T_X$.  These surfaces have no
nonidentity automorphisms by Lemma~\ref{gw:k3-center}.
The assumed rationality makes $X$ rationally connected, and the cubic
Hodge calculation gives $h^{3,1}(X)=1$
\cite[Section~2.1]{Hassett00}. Thus all hypotheses of
Proposition~\ref{gw:net-center} hold, and it gives
\begin{equation}
 \label{cub:one}
 c_u(f)=1.
\end{equation}

Both varieties are Mori fibre spaces over a point.  Indeed, they are
smooth, and integral weak Lefschetz and the exponential sequence give
$\operatorname{Pic}(X)=\Z\mathcal O_X(1)$, while adjunction gives
$-K_X=3H$ \cite[Chapter~1, Section~1.1, Lemma~1.6 and
Corollary~1.9]{HuybrechtsCubic2023}.  The ordinary Sarkisov program therefore factors $f$ into the
links used in Theorem~\ref{sf:vanishing}.  Applying that theorem to this
factorization yields $c_u(f)=0$, contradicting \eqref{cub:one}.
\end{proof}

\subsection{The classical and categorical inputs}

We use the following established results in their untwisted form, with
the labellings and admissible discriminants defined in the introduction.
Hassett's period and lattice theorems
\cite[Theorems~1.0.1, 1.0.2 and 5.1.3; Corollary~5.2.4]{Hassett00}
give the nonempty irreducible divisor $\mathcal C_d$ when
$d>6$ and $d\equiv0,2\pmod6$.  For admissible $d$, a labelled cubic has
an associated primitively polarized K3 surface $(S,L)$ of degree $d$,
with a Hodge isometry
\begin{equation}
 \label{cub:associated}
 K^\perp \simeq L^\perp(-1),
\end{equation}
where the surface cup pairing is reversed. Hassett's marked space also
records an identification of the distinguished rank-two lattice $K$
with a fixed abstract lattice, preserving $h^2$. The map forgetting
this identification is finite \cite[Section~5.2 and
Proposition~5.2.1]{Hassett00}. His marked space admits an open immersion
into the moduli space of degree-$d$ polarized K3 surfaces
\cite[Corollary~5.2.4]{Hassett00}.

For such a cubic, there exists a smooth projective K3 surface $S'$ and
an exact $\C$-linear equivalence
\begin{equation}
 \label{cub:categorical}
 \operatorname{Ku}(X)\simeq D^b(\operatorname{Coh}(S')).
\end{equation}
This is \cite[Corollary~29.7]{BLMPNS21}; it applies to every cubic having
a Hodge-theoretically associated K3 surface.  Its target is the ordinary
derived category.  Already
\cite[Theorem~1.1]{AddingtonThomas14} gives \eqref{cub:categorical} on a
Zariski open dense subset of each admissible $\mathcal C_d$, which would
suffice below.  Those equivalences are Fourier--Mukai equivalences, so
they have the exactness and $\C$-linearity in Theorem~\ref{thm:main}.

\subsection{The very general transcendental lattice}

\begin{lemma}
\label{cub:generic}
For each admissible $d$, a very general point $X\in\mathcal C_d$
satisfies
\begin{equation}
 \label{cub:generic-data}
 \operatorname{rank}A_X=2,\quad
 \operatorname{rank}T_X=21,\quad
 |\det T_X|=d,\quad
 \operatorname{End}_{\mathrm{Hdg}}(T_X\otimes\Q)=\Q.
\end{equation}
Moreover, $T_X\otimes\Q$ is simple.  All these assertions hold
simultaneously with \eqref{cub:categorical} outside a countable union
of proper closed algebraic subsets of $\mathcal C_d$.
\end{lemma}

\begin{proof}
We spell out the exceptional sets, including their algebraicity.
Pass to Hassett's marked space and then add finite level structure.
Fix a smooth irreducible component that is finite and surjective over
a dense open $U\subset\mathcal C_d$, and denote it by $U'$.
The rank-two marking makes $K$ constant. Its orthogonal
$\mathbb T=K^\perp$ is an integral polarized variation on $U'$;
we trivialize this complementary local system only on period charts.
All exceptional loci below are constructed on this one fixed cover.

The cubic Hodge numbers and middle lattice are
\[
 h^{3,1}=h^{1,3}=1,\quad h^{2,2}=21,\quad b_4=23,
 \qquad
 H^4(X,\Z)\simeq \langle1\rangle^{\oplus21}
                         \oplus\langle-1\rangle^{\oplus2};
\]
see \cite[Section~2.1 and Proposition~2.1.2]{Hassett00}.
Thus $\mathbb T$ has rank $21$ and signature $(19,2)$ for the cubic
cup pairing.  In a marked period chart its period line varies in an
open subset of
\[
 \mathcal D_T=
 \left\{[\omega]\in\mathbb P(T\otimes\C):
      (\omega,\omega)=0,\ (\omega,\overline\omega)<0\right\}.
\]
Here $T$ denotes the underlying fixed lattice in that chart.  The
openness follows equivalently from the associated K3 period description
\eqref{cub:associated}. The domain $\mathcal D_T$ is analytically open in a smooth
irreducible projective quadric, and this quadric spans
$\mathbb P(T\otimes\C)$.

For $0\ne t\in T\otimes\Q$, consider the locus in $U'$ where some
flat translate of $t$ is a Hodge class. After clearing denominators
and Tate twisting, \cite[Theorem~1.1 and Corollary~1.3]{CDK95} makes
this locus closed algebraic, with finitely many local branches.
Here one fixed denominator works for the whole monodromy orbit, and
the flat polarization has the same value on all its translates.
The cited corollary obtains this orbit locus from images of connected
components of the bounded-norm space of Hodge classes, which the
theorem makes finite over the base.
In a period chart each branch lies in
$(t')^\perp\cap\mathcal D_T$ for a nonzero flat translate $t'$.
Each such section is proper because the period domain spans
$T\otimes\C$; hence the global locus is proper as well.
There are countably many rational $t$. Since
$H^4(X,\Q)=K_\Q\oplus K_\Q^\perp$, excluding these loci removes every
Hodge class outside $K_\Q$. Thus $A_X\otimes\Q=K\otimes\Q$;
saturation gives $A_X=K$ and consequently $\operatorname{rank}T_X=21$.

The middle lattice is unimodular, and both $A_X$ and $T_X$ are
primitive.  To recall the determinant argument, write $L=H^4(X,\Z)$.
Every homomorphism $A_X\to\Z$ extends to $L$, since $L/A_X$ is
torsion-free.  Unimodularity therefore makes the restriction map
$L\to A_X^*$ surjective with kernel $T_X$.  Reducing modulo $A_X$
gives
\[
 L/(A_X\oplus T_X)\simeq A_X^*/A_X.
\]
Interchanging $A_X$ and $T_X$ gives the same quotient as $T_X^*/T_X$.
The two discriminant groups have the same order; hence
$|\det T_X|=\det A_X=d$.

Next fix a nonscalar $a\in\operatorname{End}_{\Q}(T\otimes\Q)$.
If $a$ is a Hodge endomorphism, then $\C\omega$ is an eigenline for
$a$.  Hence its Hodge locus is contained in
\[
 \bigcup_{\lambda}\mathbb P\ker(a-\lambda)\cap\mathcal D_T.
\]
This is a finite union of proper linear sections of the period quadric.
Indeed, an eigenspace equal to $T\otimes\C$ would make $a$ scalar;
and an irreducible quadric spanning the ambient vector space cannot
be contained in a finite union of proper linear subspaces.
There are countably many rational endomorphisms $a$. Apply the same
argument to the weight-zero polarized variation
$\operatorname{End}(\mathbb T)$, clearing denominators. The locus where
some flat translate of $a$ is Hodge is closed algebraic with finitely
many local branches. Monodromy acts by conjugation and preserves
nonscalarity, so the eigenline argument makes every local branch
proper. These global loci are therefore proper. Excluding them gives
$\operatorname{End}_{\mathrm{Hdg}}(T_X\otimes\Q)=\Q$.

For completeness, simplicity also follows directly from the Hodge
numbers.  By polarizability, a rational Hodge substructure has a Hodge
complement.  In any nontrivial direct-sum decomposition of $T_X\otimes\Q$,
only one summand could contain the one-dimensional $(3,1)$ part, and
that summand would also contain its conjugate.  Every other summand
would consist entirely of rational $(2,2)$ classes, contrary to the
definition of $T_X$.

Finally take the finite images in $U$ of the exceptional loci in $U'$.
These images are closed and proper, since $U'$ is irreducible and
finite over $U$. Their closures in $\mathcal C_d$ are still proper.
Adding $\mathcal C_d\setminus U$ gives a countable union of proper
closed algebraic subsets, whose complement is nonempty over $\C$.
The categorical
assertion holds throughout $\mathcal C_d$ by
\eqref{cub:categorical}; if one uses only the Addington--Thomas theorem,
one additionally removes the complement of its dense open subset.
This proves the simultaneous assertion.
\end{proof}

\begin{proof}[Proof of Theorem~\ref{thm:main}]
Let $d_0$ be as in Proposition~\ref{cub:criterion}, and fix any admissible
$d>d_0$.  The divisor $\mathcal C_d$ is nonempty.  By
Lemma~\ref{cub:generic}, outside a countable union of proper closed
algebraic subsets of this divisor, the lattice $T_X$ has rank $21$,
absolute determinant $d$, and rational Hodge endomorphism ring $\Q$.
The same lemma supplies an exact $\C$-linear equivalence
\[
 \operatorname{Ku}(X)\simeq D^b(\operatorname{Coh}(S))
\]
for a smooth projective K3 surface $S$.  Proposition~\ref{cub:criterion}
applies to each such $X$ and proves that it is irrational.

The integer $d_0$ was fixed before choosing $d$, so the conclusion holds
very generally in every admissible $\mathcal C_d$ with $d>d_0$. Its
ineffectiveness is inherited from the uniform bound in
Theorem~\ref{sf:vanishing}. Such a cubic satisfies the categorical
condition of Kuznetsov's conjecture and fails its rationality condition;
hence the categorical-to-rational implication is false.
\end{proof}

The sequence $d=2\cdot7^j$ gives an explicit infinite subfamily of the
admissible discriminants. For $j\geq1$, one has $d>6$ and
$d\equiv2\pmod6$; the prime $2$ occurs with exponent one, $3$ does not
divide $d$, and the only odd prime divisor is $7\equiv1\pmod3$.
Thus every member satisfies \eqref{cub:admissible}, and the theorem
applies to all sufficiently large members of this sequence.

\subsection{Proofs of the consequences}\label{cub:consequences}

We finish by combining Theorem~\ref{thm:main} with the established K3
and zero-cycle results used in the introduction.

\begin{proof}[Proof of Corollary~\ref{cor:associated-k3}]
For admissible \(d\), the classical input \eqref{cub:associated} supplies
the indicated polarized K3 surface for every labelled cubic in
\(\mathcal C_d\). Theorem~\ref{thm:main} supplies irrationality on the
stated very-general set for \(d>d_0\). Since \(K\) contains \(h^2\), the
isometry gives
\[
 L^\perp(-1)\simeq K^\perp\hookrightarrow
 (h^2)^\perp=H^4(X,\Z)_{\mathrm{pr}}.
\]
The inclusion is primitive: if a nonzero integer multiple of an integral
class in \((h^2)^\perp\) is orthogonal to \(K\), then the class itself is
orthogonal to \(K\). This is the primitive polarized K3 Hodge embedding
in the cited formulation, with the stated Tate twist and pairing reversal.
\end{proof}

\begin{proof}[Proof of Corollary~\ref{cor:hilbert-square}]
For \(n\ge2\), the integer \(d=2n^2+2n+2\) is greater than six and is
congruent to zero or two modulo six. Taking \(a=1\) realizes Addington's
condition \((***)\), namely
\(d=(2n^2+2n+2)/a^2\) for integers \(n,a\) with \(a\ne0\).
That condition implies his condition \((**)\), the remaining divisibility
restrictions in \eqref{cub:admissible}; hence \(d\) is admissible.
Addington's Theorem~2 then supplies the stated birationality for every
\(X\in\mathcal C_d\), with a projective K3 surface by his convention
\cite[condition~\((***)\), Theorem~2, and the convention in the introduction]{Addington16}.
Theorem~\ref{thm:main} supplies irrationality for a very general such
\(X\) when \(d>d_0\). Finally, the values \(2n^2+2n+2\) are strictly
increasing and unbounded for \(n\ge2\), so infinitely many exceed \(d_0\).
\end{proof}

\begin{proof}[Proof of Corollary~\ref{cor:chow-diagonal}]
First let \(X\in\mathcal C_d\) be any cubic with a labelling \(K\) of
admissible discriminant \(d\). The intersection form satisfies
\((h^2,h^2)=\int_Xh^4=3\), so \(h^2\) is primitive in the integral middle
lattice: if \(h^2=m\eta\) for an integral class \(\eta\), then \(m^2\)
divides \(3\). It is therefore primitive in \(K\) and extends to a basis
\((h^2,\sigma)\) of this rank-two lattice.

Voisin's Theorem~5.6 is stated for the lattice \(P\) generated by the
independent integral Hodge classes \(h^2\) and \(\sigma\), which are
algebraic on a smooth cubic fourfold as recalled in its proof
\cite[Theorem~5.6 and its proof]{Voisin17}. Here \(P=K\), and its
discriminant in that theorem is exactly
\[
 D(\sigma)=
 \det\begin{pmatrix}
  3 & h^2\mathbin{\cdot}\sigma\\
  h^2\mathbin{\cdot}\sigma & \sigma^2
 \end{pmatrix}
 =\det K=d.
\]
Admissibility includes \(4\nmid d\). Voisin's theorem therefore gives
universal triviality of \(\operatorname{CH}_0(X)\) for every such labelled
cubic, without requiring \(K\) to be the entire lattice of Hodge classes
or \(X\) to be very general. The field-extension scope and the equivalence
with the displayed integral Chow-theoretic decomposition are given in the
introduction and equation~(1) of the same paper \cite{Voisin17}.
Theorem~\ref{thm:main} supplies irrationality for the stated very general
cubics when \(d>d_0\).
\end{proof}
